\scanpage{4}
$G$ groupe algébrique sur $\mathbb{Q}$

\[
\unsure{\mathbb{S}|\mathrm{S}}=\prod_{\mathbb{C}/\mathbb{R}}G_{m\,\mathbb{C}},\quad
\text{donc}\quad \mathbb{S}(\mathbb{R})\simeq\mathbb{C}^{*}.\quad
\text{Soit }\underset{\substack{\Vert\\
\operatorname{Ker}(\mathrm{Norme}:\mathbb{S}\to G_{m\,\mathbb{R}})}}{\mathbb{U}}\subset\mathbb{S}
\]

$\operatorname{Rep}(\mathbb{S})\approxeq$ catégorie des vectoriels réels $V_{\mathbb{R}}$, avec bigraduation sur $V_{\mathbb{C}}=V_{\mathbb{R}}\otimes_{\mathbb{R}}\mathbb{C}$,
\[
V_{\mathbb{C}}\simeq\coprod_{p,q}V_{\deleted{\unsure{p,q|\mathbb{C}}}}^{p,q}
\]
\margin{\uline{NB}\quad $\lambda\in\mathbb{C}^{*}\simeq\mathbb{S}(\mathbb{R})$ opère sur $V^{pq}$ par $\lambda^{p}\bar{\lambda}^{q}$}
satisfaisant la condition
\[
V^{p,q}=\overline{V^{q,p}}.
\]

L'hom canonique $G_{m\,\mathbb{R}}\xrightarrow{i}\mathbb{S}$ correspond à la graduation totale

\diagram{\begin{tikzcd}[ampersand replacement=\&]
\operatorname{Rep}(G) \arrow[r,"\text{oubli}"] \&
\operatorname{Modf}(\mathbb{Q}) \arrow[r,"\otimes_{\mathbb{Q}}\mathbb{R}"] \&
\operatorname{Modf}(\mathbb{R}) \arrow[r,"\otimes_{\mathbb{R}}\mathbb{C}"] \&
\operatorname{Modf}(\mathbb{C})
\end{tikzcd}}

La donnée sur le foncteur composé d'une bigraduation
\[
V_{\mathbb{C}}\simeq\coprod_{p,q}V_{\deleted{\mathbb{C}}}^{p,q}
\]
\margin{i.e. d'une préstructure de Hodge, compatible avec la structure sur $\mathbb{Q}$ donnée, et fonctorielle en l'objet de $\operatorname{Rep}(G)$}
satisfaisant
\[
V^{p,q}=\overline{V^{q,p}}
\]
équivaut à la donnée d'un hom
\[
\mathbb{S}\xrightarrow{j}G_{\mathbb{R}}
\]
Si on se donne un hom [central]
\[
i:G_m\rightharpoonupG.\;,
\]
i.e. une graduation :
\[
V=\coprod_n V^n
\]
sur le foncteur oubli \unsure{$w$|$\omega$}, alors la condition
\[
V^n_{\mathbb{C}}=\coprod_{p+q=n}V^{p,q}
\]
% archive sheet number: 3

